% !TeX root = ../main.tex
\section{Methodology}
\label{sec:methodology}

\subsection{Research Questions and Study Scope}
\label{sec:research-questions}

We organize the study around three research questions
(Table~\ref{tab:research-questions}). The computational protocol treats a user
story and its acceptance criteria as a requirement case, with the system purpose
providing context. It assesses each catalog pattern against each case
independently. The assessment distinguishes applicability, non-applicability,
and insufficient information; a positive classification requires a supporting
user-story or acceptance-criterion reference. These classifications are model
outputs, not independently established ground truth.

\begin{table*}[t]
\caption{Research questions, units of analysis, and evidence mapping.}
\label{tab:research-questions}
\centering
\small
\begin{tabular}{p{0.05\textwidth} p{0.36\textwidth} p{0.49\textwidth}}
\hline
\textbf{ID} & \textbf{Research question} & \textbf{Unit, evidence, and intended analysis} \\
\hline
RQ1 &
How does the LLM classify architectural-pattern applicability across the requirement cases, and which patterns does it identify as applicable? &
Unit: one user-story--pattern pair.
Evidence: the official exhaustive assessment log, original requirements, and pattern catalog.
Analysis: distributions of the three classifications and applicable-pattern sets by story and system, retaining zero-positive stories. \\
\hline
RQ2 &
To what extent do the patterns classified as applicable by the LLM overlap with the frozen RAG baseline's final Top-5 candidates for the same requirement cases? &
Unit: one LLM-positive user-story--pattern pair, grouped by story.
Evidence: official applicability assessments and the frozen final Top-5 rankings.
Analysis: membership-intersection counts by story and overall, with the LLM-positive set as an explicit denominator; no applicability-accuracy metric. \\
\hline
RQ3 &
How do experts assess the applicability of the LLM-selected patterns and the adequacy, grounding, consequences, and perceived usefulness of the corresponding generated ADRs? &
Unit: one expert's rating of one selected pair and its ADR on one questionnaire item.
Evidence: consolidated expert responses to Q1--Q5 and optional comments for the prepared cases.
Analysis: item-specific ordinal distributions and case-level summaries with response counts; comments contextualize ratings. Aggregation and inclusion rules must be specified before consolidated analysis. \\
\hline
\end{tabular}
\end{table*}


RQ1 characterizes the model's classification outcomes rather than measuring
their correctness. RQ2 compares the patterns classified as applicable with a
separate retrieval baseline. This baseline retrieves catalog candidates and
reranks them to a final Top-5 without making an LLM-based applicability judgment.
Overlap is therefore a descriptive relationship between two output sets, not
evidence of either approach's precision, recall, or superiority.

ADR generation is a separate stage: each pair classified as applicable supplies
the requirement case, selected pattern, and applicability evidence for one
generated ADR. Generation completeness and stored provenance establish the
availability and traceability of these records, not their architectural quality.
RQ3 examines expert assessments of these selected pairs and their corresponding
ADRs, rather than adding another automatic applicability classification.

The prepared questionnaire measures decision adequacy for the stated
requirements (Q1), decision grounding (Q2), presentation of relevant
consequences and trade-offs (Q3), perceived usefulness as a decision record
(Q4), and pattern applicability (Q5). These are five separate expert-assessed
dimensions; they are not a validated composite quality score. Optional comments
provide context for the ratings. Neither implementation feasibility, cost
savings, novelty, nor actual improvement in architectural practice is directly
measured by these items.

Experts evaluate cases for a system they know, and only pairs selected by the
LLM are presented. Consequently, RQ3 can characterize expert judgments of the
presented recommendations but cannot identify missed applicable patterns or
estimate recall over the full set of assessed pairs. Expert agreement with a
recommendation does not by itself establish implementation success or objective
correctness. Repeated ratings from an expert and multiple cases from a system
must not be treated as independent replications when interpreting the results.
Conclusions remain bounded to the sampled requirements, catalog, generated
records, and participating experts.

% Claim-to-evidence matrix for Task 1.1. The E1--E6 paths below are relative to
% docs/20261012.1000-PaperInitialVersion/, not this relocated LaTeX directory.
% These comments are provenance pointers, not external LaTeX build inputs.
% E1: ../20261012.0933-InitialArtifacts/adr4ai_app/experiments/pattern_selection/
%     assess_pattern.py and outputs/assessments.jsonl.
% E2: ../20261012.0933-InitialArtifacts/adr4ai_app/data/user_histories_orig.csv
%     and catalogs/patterns/.
% E3: ../20261012.0933-InitialArtifacts/adr4ai_app/experiments/rag/
%     README.md and outputs/rag_baseline.csv (not reproduction or pilot outputs).
% E4: ../20261012.0933-InitialArtifacts/adr4ai_app/research_results/
%     adr_generation_results.csv and compile_adr_results.py;
%     ../20261012.0933-InitialArtifacts/adr4ai_app/experiments/adr_generation/
%     generate_adr.py and outputs/*.json (excluding attempts/).
% E5: ../20261012.0933-InitialArtifacts/adr4ai_app/research_results/google_forms/
%     create_expert_evaluation_form.gs, EVALUATION_QUESTIONS, and README.md;
%     research_results/expert_evaluation_cases_pt.csv.
% E6: Future frozen/consolidated expert-response dataset and documented analysis.
%     Researcher reports preliminary responses; no response data inspected here.
%
% C1 / RQ1: generated label distributions and positive sets -> E1 + E2.
%     Allowed: descriptive counts/sets. Excluded: expert-confirmed correctness.
% C2 / RQ2: final Top-5 membership overlap -> E1 + E3, joined by story/pattern ID.
%     Allowed: intersection counts. Excluded: precision, recall, superiority.
% C3 / generation: one traceable ADR per selected pair -> E4 plus E1.
%     Allowed: completeness/provenance once verified. Excluded: semantic quality.
% C4 / RQ3 constructs: Q1--Q4 ADR dimensions; Q5 pattern applicability -> E5.
%     Allowed now: instrument description. Ratings/results require E6.
% C5 / RQ3 findings: observed expert ratings/comments -> E5 + E6.
%     Blocked pending consolidation; no synthetic or favorable mock findings.
%     Excluded: missed-pattern recall, population-wide effectiveness, deployment
%     success, cost savings, and a composite score without justification.

\subsection{Industrial Requirement Sample}
\label{sec:requirement-sample}

The study uses a bounded industrial sample comprising
\StudySystemCount{} systems, \StudyStoryCount{} user stories
(\StudyStoriesPerSystem{} per system), and \StudyCriterionCount{} acceptance
criteria (Table~\ref{tab:dataset-summary}). The supplied data-scope document
defines this as a minimal, delimited sample for studying architectural-pattern
selection and decision documentation for AI-enabled systems. It is not a
representative sample of industrial requirements or AI-enabled systems.
The available artifacts do not establish a probabilistic sampling procedure;
we therefore do not interpret these counts as population estimates.

% !TeX root = ../main.tex
% Generated by processing/scripts/prepare_dataset_summary.py.
\begin{table}[t]
\caption{The bounded industrial requirement sample.}
\label{tab:dataset-summary}
\centering
\begin{tabular}{lrr}
\hline
\textbf{System ID} & \textbf{User stories} & \textbf{Acceptance criteria} \\
\hline
SYS01 & 3 & 6 \\
SYS02 & 3 & 11 \\
SYS03 & 3 & 7 \\
SYS04 & 3 & 11 \\
\hline
Total & \StudyStoryCount{} & \StudyCriterionCount{} \\
\hline
\end{tabular}
\end{table}


Each requirement case consists of the system purpose, one stakeholder user
story, and all of that story's acceptance criteria. The source CSV stores one
acceptance criterion per row, repeating the corresponding purpose and story.
Grouping by the stable user-story identifier reconstructs the case without
rewriting the original Portuguese requirement text or changing criterion order.
The criterion-row count must not be confused with the number of requirement
cases.

System purpose provides contextual information, but is not sufficient on its
own to establish pattern applicability. Positive evidence must be grounded in
the user story and/or its acceptance criteria. The model is not authorized to
infer unstated technologies, requirements, or constraints from the purpose.

The source uses generic names \texttt{System\_01}--\texttt{System\_04}; the
paper uses the corresponding stable identifiers \texttt{SYS01}--\texttt{SYS04}.
Story and criterion identifiers preserve traceability within each system
without exposing real system names. The organization is not named in this
review draft. According to the supplied scope document, the study input is
restricted to purposes, identifiers, user stories, and acceptance criteria;
it does not include operational databases, personal data, credentials, customer
records, or production content. This delimitation does not itself establish
participant consent, ethics approval, or permission for unrestricted public
release.

% Task 1.2 provenance: original data/user_histories_orig.csv and the supplied
% Escopo_Dados_Estudo_Cagece_v4.pdf in the initial-artifact collection.
% Generated counts/table and audit: processing/scripts/prepare_dataset_summary.py
% and
% processing/outputs/dataset-summary-audit.json, relative to the project root.

\subsection{Architectural Pattern Catalog}
\label{sec:pattern-catalog}

We use \StudyPatternCount{} pattern records from the repository associated
with Heiland et al.'s multivocal literature review of design patterns for
AI-based systems~\cite{heiland2023patterns,aiPatternsRepository}.
For provenance verification, we compared the local files with the public
repository at revision
\texttt{9ff3e61} (the full revision is recorded in the artifact audit).
All \StudyPatternCount{} records matched the corresponding repository files
byte-for-byte, as checked using Git blob hashes. This identifies the source
snapshot of the catalog; it does not establish the effectiveness of the
individual patterns.

Each pattern is stored as one JSON record with a stable three-digit identifier,
from \texttt{001} through \texttt{070}. The applicability prompt supplies exactly
five catalog fields: \emph{ID}, identifying the record; \emph{name}, naming
the pattern; \emph{motivation}, describing the problem or need;
\emph{solution}, describing its proposed response; and \emph{consequences},
describing resulting effects and trade-offs. Identifiers remain strings,
preserving leading zeros and enabling joins with the assessment and ADR
artifacts. Source text is preserved rather than normalized or rewritten.

The records also contain aliases, examples, related-pattern identifiers,
categories, and resource identifiers. These additional fields are not supplied
to the applicability model. Resource identifiers provide source pointers within
the repository, not independently verified bibliographic descriptions in the
local catalog. We therefore attribute the collection to the verified repository
and its associated publication rather than infer individual source citations
from numeric resource codes.

All five input fields are present as strings in every record. Explicitly empty
motivation or consequences strings are retained as empty inputs, without
inventing replacement descriptions. A missing field or non-string value is an
input-validation failure, not an invitation to repair the source. This
preservation policy keeps the assessed catalog traceable while acknowledging
uneven descriptive completeness.

The catalog includes general architectural patterns, such as Client--Server,
alongside patterns concerning machine-learning safeguards and model serving.
We use the records as a bounded candidate set, not as equally mature,
mutually independent, or universally applicable solutions. Catalog inclusion
alone does not establish applicability to a requirement case; that relationship
is assessed separately.

% Task 1.3 provenance: processing/outputs/catalog-summary-audit.json and
% processing/outputs/catalog-upstream-verification.json (project-root paths).
% Bibliographic metadata checked on 2026-10-12 against Crossref's DOI record
% 10.1109/CAIN58948.2023.00035 and the pinned SWE4AI/ai-patterns repository.

\subsection{Exhaustive Pattern-Applicability Assessment}
\label{sec:applicability-experiment}

The experiment crosses all \StudyStoryCount{} requirement cases with all
\StudyPatternCount{} catalog patterns, yielding
$\StudyStoryCount \times \StudyPatternCount = \StudyAssessmentCount$ planned
pairwise assessments. Each request contains one complete requirement case and
one pattern, with no prior assessment or conversation history supplied. This
exhaustive design does not retrieve or prefilter candidates. Separate requests
do not imply statistical independence: assessments share stories, systems,
catalog entries, and a model. The resulting classifications are model
predictions, not expert-provided ground-truth labels.

\paragraph{Model and structured response}
The recorded model identifier is \texttt{\ApplicabilityModel}, with reasoning
effort \texttt{\ApplicabilityReasoningEffort}. The implementation uses the
Responses API with strict JSON-schema output named
\texttt{pattern\_assessment}. The schema requires exactly three fields:
\texttt{classification} (one of the three labels below),
\texttt{requirement\_refs} (an array of strings), and \texttt{rationale}
(a string); additional properties are disallowed. Each stored result also
includes the story ID, pattern ID, model identifier, and reasoning effort.
No temperature or seed is set in the request, and the recorded model identifier
does not pin a dated model snapshot.

\paragraph{Decision and grounding rules}
The instructions distinguish:
\begin{itemize}
  \item \texttt{applicable}: the story and/or its acceptance criteria provide
  positive evidence of a need, problem, constraint, or quality concern addressed
  by the pattern.
  \item \texttt{not\_applicable}: those requirements provide evidence that the
  pattern does not fit the stated need, problem, constraint, or context.
  \item \texttt{insufficient\_information}: the story and criteria do not provide
  enough evidence to establish either applicability or non-applicability.
\end{itemize}
System purpose supplies context but cannot alone justify applicability.
The prompt prohibits assuming unstated requirements, technologies, constraints,
or system characteristics. Absence of contradiction, possible usefulness, or
the possibility of implementing, testing, or improving a requirement using a
pattern is not sufficient positive evidence. Assessment concerns underlying
architectural intent rather than literal terminology. An applicable response
must cite at least one identifier belonging to the current story or one of its
criteria, and briefly explain the evidence--pattern relationship; it must not
propose an architecture, implementation, technology, or new requirement.
These are prompting and reference-validation constraints, not proof that a
model's rationale is semantically correct.

\paragraph{Execution, retries, and validation}
Execution is sequential, in CSV story order and ascending pattern-ID order.
SDK retries are disabled. The runner permits at most
\ApplicabilityMaxAttempts{} attempts per pair, revisiting unresolved pairs
in up to \ApplicabilityMaxAttempts{} ordered passes rather than immediately
retrying each failed request. Before a request, it atomically persists an
attempt reservation, so interrupted requests consume the attempt budget.
A request failure updates that entry with its error type and a whitespace-
normalized message, truncated to 300 characters with the configured API key
redacted. Successful validation removes the pair from this unresolved-failure
state.

The response must have completed, contain no refusal, and provide parseable
structured output. Local validation checks the exact fields, label, reference
array and identifier membership, a nonempty reference array for
\texttt{applicable}, and a string-valued rationale. These checks do not require
a nonempty rationale string or test the factual entailment of cited text.
Only validated results are appended to the official log, flushed, and synced
to disk; storage failures stop execution rather than risk duplicate writes.
A nonblocking file lock prevents concurrent official runners.

On resume, the runner validates existing JSONL records, model configuration,
pair uniqueness, and failure-state attempt counts. Malformed records,
unknown identifiers, and interrupted writes lacking a final newline stop
the run rather than being repaired or discarded. Completed pairs are skipped;
unresolved attempt counts are retained, and exhausted pairs are not retried.
The final on-disk check requires exactly all \StudyAssessmentCount{} expected
pairs and \StudyPatternCount{} records per story before exporting
story-specific applicable-pattern selections.

\paragraph{Official dataset and replication artifact}
The official dataset is the top-level \texttt{assessments.jsonl} under the
pattern-selection experiment's \texttt{outputs} directory, accompanied by
\texttt{failures.jsonl}. Separate \path{pilot_SYS01-US01} directories
(including the \texttt{\_v1} and \texttt{\_v2} variants) and
\path{test_assessment.json} and \path{test_assessment_v2.json} are excluded,
not merged with official records. An offline audit verifies complete official
coverage and records input/source hashes, the exact instructions, JSON schema,
and input template. The template presents system purpose, identified story and
criteria, followed by the five pattern fields in Section~\ref{sec:pattern-catalog};
the full instructions remain in the artifact rather than this space-bounded
description.

The frozen log contains \StudyAssessmentCount{} validated pair outcomes and
an empty final unresolved-failure file. This establishes completion, not that
no retries occurred: successful records omit attempt counts, and successful
retries remove failure-state entries. Thus, \StudyAssessmentCount{} describes
the planned first-pass requests and retained outcomes, not an independently
verified total number of API attempts. We do not claim exact stochastic
reproduction from the stored configuration.

% Task 1.4 provenance: original pattern_selection/{assess_pattern,
% run_experiment,prepare_data,run_pilot}.py and top-level outputs/*.jsonl.
% Offline audit/prompt/schema: processing/outputs/applicability-protocol-audit.json.
% Reproduction/check: processing/scripts/prepare_applicability_protocol.py.
